\section{Dimension four: the cubic model}\label{sec:cub}

Here $n=2$ and the target is $\Bo(1)$. We construct the normal disc-bundle model associated with a smooth cubic
in $\CP^2$, as in Biran's decomposition \cite{Biran2001}, and a transfer of its truncated regions to the ball.
The bundle has degree $9$ over the cubic. The base is a fixed hexagonal torus with an explicit periodic area density
of integral three. The model is written on its cover $\R^2\times\C$, with two explicit deck maps.
Its horizontal coefficient is $1-3p$, so the fibre interval stops below $p=1/3$; this gives the condition $3c_i<C$.
The integrated profile tends to $1/2$ as the height tends to $1/3$.
An explicit unit lift satisfying the Fermat equation supplies the transport and return. A cube-root symmetry,
together with injectivity modulo phase, ensures that the final ball chart is defined.

\subsection{The model}

Let $\zeta=e^{2\pi i/3}$ and let $\Gamma=(\R/\Z)^2$, identified with
$\C/(\Z+\zeta\Z)$ by $(u,v)\mapsto u+\zeta v$.
The theta formulas in \S\ref{sec:fer-torus} give a smooth map $Z:\R^2\to\C^3$ with
\[
 |Z|=1,\quad Z_0^3+Z_1^3+Z_2^3=0,\quad
 Z(u+1,v)=-Z(u,v),\quad Z(u,v+1)=\phi(u,v)Z(u,v),\quad |\phi|=1.
\]
Representatives related by $(Z,U)\mapsto(e^{i\chi}Z,e^{3i\chi}U)$ are called gauge equivalent.
Forms are evaluated on the cover and on horizontal unit representatives.

\begin{definition}[The cubic cover model]\label{def:cub-model}
\lean{SymplecticPacking.Return.Ncoord,SymplecticPacking.Bundle.Omega0loc,SymplecticPacking.Embed.deckU,SymplecticPacking.Embed.deckV,
SymplecticPacking.Bundle.pE}
Use the deck maps
\[
 T_u(u,v,w)=(u+1,v,-w),\qquad T_v(u,v,w)=(u,v+1,\phi(u,v)^3w)
\]
on $\widetilde N=\R^2\times\C$; $N$ denotes these cover data with their deck identification.
Put
\[
 p=\pi|w|^2,\quad \eta=\frac{1}{2\pi}\operatorname{Im}\sum_j\overline{Z_j(u,v)}\,dZ_j(u,v),\quad
 \sigma=d\eta,\quad \alpha=\frac{d\arg w}{2\pi}-3\eta,
\]
\[
 \Omega_0=\sigma+d(p\alpha)
 =\omega_V+(1-3p)\sigma-3dp\wedge\eta,\quad
 \Delta=[0,h],\quad 0<h<1/3,\quad \Hb(p)=3(1-3p).
\]
The expression $p\,d\arg w/(2\pi)=\frac12\operatorname{Im}(\bar w\,dw)$ makes the form smooth at $w=0$.
The marked set $\mathcal Z$ consists of the three torus points at which $Z_0=0$; its inverse image in the cover is
$\widetilde{\mathcal Z}=\mathcal Z+\Z^2$, using the three representatives given below.
The cover-to-representative map is $\mathsf N(u,v,w)=(Z(u,v),w)$.
\end{definition}

The deck transformations preserve $p$ and $\Omega_0$ and map $\mathsf N$ to gauge-equivalent pairs.
The density $\sigma=\sigma_\Gamma(u,v)du\wedge dv$ is smooth, periodic and strictly positive, with
\[
 \int_0^1\int_0^1\sigma_\Gamma(u,v)\,du\,dv=3
\]
(Lemma~\ref{lem:fer-sigma}). In particular $\Omega_0$ is nondegenerate for $p<1/3$.
The construction uses these cover identities directly.

\begin{lemma}[Integrated profile]\label{lem:cub-volume}
For $0<h<1/3$,
\[
 \int_0^h\Hb(p)\,dp=3h-\frac92h^2,\qquad
 \frac12-\int_0^h\Hb(p)\,dp=\frac12(1-3h)^2.
\]
Thus this scalar budget tends to $1/2$ as $h\uparrow1/3$.
\end{lemma}
\begin{proof}
Integrate the affine function $3-9p$ and complete the square.
\end{proof}


\subsection{Handle data}

\emph{Torus coordinates.} Let $\zeta=e^{2\pi i/3}$ and $\Lambda_0=\Z+\zeta\Z$. The Weierstrass parametrisation
\[
 \Phi_\Gamma(z)=\big[\,2c_F\,\wp(z):-(1+d_F\,\wp'(z)):-(1-d_F\,\wp'(z))\,\big],\qquad z=u+\zeta v,\ (u,v)\in\R^2,
\]
with $\Phi_\Gamma(0)=[0:1:-1]$, $\wp=\wp_{\Lambda_0}$, $g_3=g_3(\Lambda_0)\neq0$, $c_F=g_3^{-1/3}$ and
$d_F=(3g_3)^{-1/2}$, is injective modulo the lattice as a map to gauge classes of Fermat vectors. An explicit theta lift $Z(u,v)$ without poles makes $\sigma$ a smooth form on $(\R/\Z)^2$ and $Z^*\eta_H$ a smooth form on
$\R^2$ with explicit lattice periods; the base data below are smooth (\S\ref{sec:fer}). In these
coordinates $\nu=du\wedge dv$ is a positive area form with $\int\nu=1$, and
\[
 \mathcal Z=\{b_1,b_2,b_3\},\qquad b_1=(0,0),\quad b_2=(\tfrac13,\tfrac23),\quad b_3=(\tfrac23,\tfrac13).
\]
(The point $b_1$ is $\Phi_\Gamma(0)=[0:1:-1]$; $b_2$, $b_3$ are the two zeros of $\wp$, which are the non-lattice
fixed points of $z\mapsto\zeta z$.)

\begin{proposition}[Handle data on the cubic]\label{prop:cub-handle}
\lean{SymplecticPacking.Bundle.torusPoint, SymplecticPacking.Bundle.bpoints, SymplecticPacking.Bundle.Z0_eq_zero_iff, SymplecticPacking.Bundle.annulusPlus, SymplecticPacking.Bundle.annulus, SymplecticPacking.Bundle.bB, SymplecticPacking.Bundle.betaB, SymplecticPacking.Bundle.stripB, SymplecticPacking.Bundle.zetaJ, SymplecticPacking.Bundle.supports, SymplecticPacking.Bundle.Pprim, SymplecticPacking.Bundle.Pprim_periodic, SymplecticPacking.Bundle.d_Pprim, SymplecticPacking.Bundle.frameA, SymplecticPacking.Bundle.lambdaC_descends, SymplecticPacking.Bundle.d_lambdaC, SymplecticPacking.Bundle.angularForm_frameA}
\uses{def:cub-model, M-def:hdl-handle-data, prop:fer-handle, lem:fer-frame}
In the torus coordinates $(u,v)$ the following are handle data for the cubic model
$(\Gamma,N,[0,h],\Omega_0,\mathcal Z)$ (\Mref{Definition}{def:hdl-handle-data}):
\begin{enumerate}[label=(\alph*)]
\item the annuli $\mathcal A=\{\frac18\le u\le\frac5{24}\}\subset\mathcal A^+=\{\frac1{12}\le u\le\frac14\}$ with handle
coordinates $(s,v)=(12u-\frac32,\,v)$, so that $\mathcal A=[0,1]\times\R/\Z$ and $\mathcal A^+=[-\frac12,\frac32]\times\R/\Z$;
\item the dual band $\mathcal B=\{\frac19<v<\frac29\}$, containing the closed clock band $[\frac18,\frac5{24}]$; its active intersection with $\mathcal A^+$ is $[\frac1{12},\frac14]\times[\frac18,\frac5{24}]$, with $v_-=\frac18$, $v_+=\frac5{24}$;
\item the closed one-form $\beta=b(v)\,dv$, $b=\varphi/\int_0^1\varphi$, where
$\varphi(v)=a_{\rm flat}(v-\frac18)\,a_{\rm flat}(\frac5{24}-v)$ on $[0,1]$, extended $1$-periodically, with
$a_{\rm flat}(x)=e^{-1/x}$ for $x>0$ and $a_{\rm flat}(x)=0$ for $x\le0$;
\item the area form $\nu=du\wedge dv=\frac1{12}\,ds\wedge dv$, so $c_\nu=\frac1{12}$;
\item the primitive operator $P_\Gamma$: for a smooth function $\mathfrak b$ on $(\R/\Z)^2$ with $\int\mathfrak b\,\nu=0$,
\[
\begin{aligned}
 P_\Gamma(\mathfrak b\,\nu)&=A_0\,dv-B_0\,du,\\
 A_0(u,v)&=\int_0^u\big(\mathfrak b(s',v)-\mathfrak b_0(v)\big)ds',\\
 \mathfrak b_0(v)&=\int_0^1\mathfrak b(s',v)\,ds',\qquad B_0(v)=\int_0^v\mathfrak b_0(t)\,dt;
\end{aligned}
\]
\item the unitary frame $(Z_0/|Z_0|)^3$ of $N$ over $\mathcal A^+$, with fibre coordinate $w_{\mathcal A}$ defined by
$U=(Z_0/|Z_0|)^3w_{\mathcal A}$, angle $\theta_{\mathcal A}=\arg w_{\mathcal A}/2\pi$, and the toric one-form
\[
 \Gamma_{\mathcal A}(p)=(1-3p)\,\lambda_\Gamma,\qquad \lambda_\Gamma=\eta_H-\frac1{2\pi}\,d\arg Z_0 ,
\]
so that $\Omega_0=\omega_V+d\Gamma_{\mathcal A}$ on $N|_{\mathcal A^+}$, with $\omega_V=dp\wedge d\theta_{\mathcal A}$;
\item the coordinate discs $U_j$ of radius $\frac1{100}$ about $b_j$ in the centred coordinates
$(u-u(b_j),v-v(b_j))$.
\end{enumerate}
\end{proposition}

\begin{proof}
(a), (b) The maps $u\mapsto12u-\frac32$ and $v\mapsto v$ identify the closed bands as stated, and $ds\wedge dv=12\,du\wedge dv$ is
positive. Modulo $1$, the $u$-coordinates $0,\frac13,\frac23$ of the $b_j$ have distance at least $\frac1{12}$ from
$[\frac1{12},\frac14]$, and their $v$-coordinates $0,\frac23,\frac13$ have positive distance from $[\frac19,\frac29]$;
so $\mathcal A^+$ and $\mathcal B$ avoid $\mathcal Z$. The base is closed, so there is no boundary collar.
(c) $\varphi$ is smooth on $\R$ because $a_{\rm flat}$ is flat at $0$; it vanishes outside $[\frac18,\frac5{24}]\subset(0,1)$ and is
positive exactly on $(\frac18,\frac5{24})$. Hence $b\ge0$ is smooth and $1$-periodic, $b>0$ exactly on $(v_-,v_+)$,
$\int_0^1b=1$, $\operatorname{supp}\beta\subset\mathcal B$, $\beta$ is closed because it depends on $v$ only, and
$\int_{\{s\}\times\R/\Z}\beta=1$.
(d) is the positivity of $\nu$ (\S\ref{sec:fer}).
(e) $A_0$ and $B_0$ are $1$-periodic because $\int_0^1(\mathfrak b-\mathfrak b_0)\,ds'=0$ and $\int_0^1\mathfrak b_0=\int\mathfrak b\,\nu=0$,
and $d(A_0\,dv-B_0\,du)=(\mathfrak b-\mathfrak b_0)\,du\wedge dv+\mathfrak b_0\,du\wedge dv=\mathfrak b\,\nu$; the operator is linear and
preserves smooth dependence on parameters, with all jets, by differentiation under the integral sign. The support
condition of \Mref{Definition}{def:hdl-handle-data}(e) is void on a closed base.
(f) $\mathcal Z=\{b_1,b_2,b_3\}$ is disjoint from $\mathcal A^+$ by (a), so $Z_0\ne0$ on $\mathcal A^+$ and the frame
is defined there. The identity $\Omega_0=\omega_V+d\Gamma_{\mathcal A}$ is Lemma~\ref{lem:fer-frame}: $\lambda_\Gamma$ is basic
with $d\lambda_\Gamma=\sigma$, and $\arg U=3\arg Z_0+\arg w_{\mathcal A}$ gives $\alpha=d\theta_{\mathcal A}-3\lambda_\Gamma$. The form
$\Gamma_{\mathcal A}$ is a smooth family of base one-forms depending on $p$, hence toric.
(g) The $b_j$ have mutual distance at least $\frac13$ in the flat metric of the torus, and by (a), (b) distance at least
$\frac1{12}$ from $\mathcal A^+\cup\mathcal B$; both bounds exceed $\frac2{100}$. So the closed discs $\overline U_j$ are disjoint from
each other and from $\mathcal A^+\cup\mathcal B$.
\end{proof}

The base is closed, so there is no collar ($\chi_\partial=0$) and no finite atlas: the primitive operator is the global
operator $P_\Gamma$ of (e), and the correction region of \Mref{Definition}{def:hdl-handle-data}(d) is the whole torus, on
which $\kappa(p)=(1-3p)\sigma\ge(1-3h)\sigma$ for $p\le h$. The background is $e_h\nu$ (Definition~\ref{def:fer-kappa}). The
twist, the three paths, the Moser map and the layer maps over these data are constructed in \S\ref{sec:shd-four}.

\subsection{Transfer into the ball}

The transfer is constructed by a deformation to the normal cone, following the scheme of \cite{OpenAI2026}.
It embeds each working region $\{p<h\}$, away from the marked fibres and modulo deck identification, into the
ball. Let $\mathcal F(Z)=Z_0^3+Z_1^3+Z_2^3$ and consider the family
$X_T=\{\mathcal F(Z)=TW\}$, where $W$ has weight three. For $T=1$ the equation says $W=\mathcal F(Z)$, so $X_1$ is a graph
over $\CP^2$. For $T=0$ it says $\mathcal F(Z)=0$ with $W$ free, so $X_0$ is the total space of the normal bundle
$\mathcal O(3)|_E$ of the cubic $E$, which is the disc-bundle model. One Kähler potential $\Phi_E$ restricts to
all fibres $X_T$. The weight-three injection places the model in $X_0$ with its own form; the transport is the
$\Omega_E$-orthogonal lift of $\partial_T$ on a bounded region, and carries its points from
$X_0$ to $X_1$ preserving the restricted forms; projection identifies $X_1$
with $\CP^2$, where the restricted form is $\omega_{\rm FS}$ plus the exact term $d_n\,dd^c\log(1+\mathfrak u)$; and
the Moser return removes that term. Both vector fields are given in closed form, with denominators bounded below
for $0\le T\le1$ and $0\le s\le1$. To land in $\Bo(1)=\CP^2\setminus\{Z_0=0\}$ the image must avoid the line
$\{Z_0=0\}$. Every step commutes with the symmetry $\mathsf r\colon Z_0\mapsto\zeta Z_0$, whose fixed points in $\CP^2$ are that line
and the point $[1:0:0]$. A map that commutes with $\mathsf r$ and is injective modulo gauge sends non-fixed points to
non-fixed points, so the model over $E\setminus\{Z_0=0\}$ is carried off the line
(Proposition~\ref{prop:rtc-avoid}).

\begin{construction}[The cubic transfer]\label{cons:cub-transfer}
\uses{def:cub-model, cons:rtc-injection}
Fix $h<d_n<1/3$. Use pairs $(Z,W)$ with $|Z|=1$, with the gauge relation
$(Z,W)\sim(e^{i\chi}Z,e^{3i\chi}W)$.
The ambient form is defined on horizontal representatives by the potential
\[
 \Phi_E(Z,W)=\log|Z|^2+d_n\log(1+|W|^2/|Z|^6),\qquad \Omega_E=dd^c\Phi_E.
\]
The constraint is $X_T=\{(Z,W):|Z|=1,\ \mathcal F(Z)=TW\}$.
In unit representatives put $p_W=d_n|W|^2/(1+|W|^2)$, $A=1-3p_W$ and $B=d_n/(1+|W|^2)^2$. For $\|Z\|=1$ put
$f=\mathcal F(Z)$, $\mathbf n=3\bar Z^{2}-3\bar fZ$ (with $\bar Z^2=(\bar Z_0^2,\bar Z_1^2,\bar Z_2^2)$), $G=\|\mathbf n\|^2$,
$N_{\rm den}=G/A+T^2/B$ and $\mathfrak u=|\mathcal F(Z)|^2/\|Z\|^6$. The transfer is the composite of the following four maps.
\begin{enumerate}[label=(\arabic*)]
\item The \emph{weight-three injection} $\iota[Z,U]=[Z,\sqrt\pi\,U/\sqrt{d_n-p}]$ for unit pairs with $p<d_n$.
\item The \emph{transport} $\tau_1$, the time-one map, from $T=0$ to $T=1$, of the field
\[
 \dot Z=\frac{W\,\mathbf n}{A\,N_{\rm den}},\qquad \dot W=-\frac{T\,W}{B\,N_{\rm den}},\qquad \dot T=1
\]
on unit representatives in the gauge $Z^\dagger\dot Z=0$.
\item The projection $\mathrm{pr}(Z,W)=Z$ from the graph constraint at $T=1$.
\item The \emph{Moser return} $\psi_1^{-1}$, where $\psi_s$ is the unit-representative flow of
\[
 Y_s=-\frac{d_n}{2}\,\frac{f\,\mathbf n}{(1+\mathfrak u)(A_s+B_sG)},\qquad A_s=1-3sd_n\frac{\mathfrak u}{1+\mathfrak u},\quad
 B_s=\frac{sd_n}{(1+\mathfrak u)^2},
\]
with the graph forms $\omega^{\rm gr}_s=\omega_{\rm FS}+s\,d_n\,dd^c\log(1+\mathfrak u)$, $0\le s\le1$.
\end{enumerate}
Finally set $\mathsf r Z=(\zeta Z_0,Z_1,Z_2)$ and $\mathsf r(Z,W)=(\mathsf r Z,W)$. This preserves the central constraint and the gauge relation.
\end{construction}

The ball chart $b_{\rm ball}$ (Lemma~\ref{lem:rtc-ball-chart}) is applied to the output of $\psi_1^{-1}$. The resulting
cover transfer $\mathcal T=b_{\rm ball}\circ\psi_1^{-1}\circ\mathrm{pr}\circ\tau_1\circ\iota\circ\mathsf N$ is
defined on $\{p<h\}$ off the marked fibres, pulls $\omega_0$ back to $\Omega_0$ and identifies exactly the deck-equivalent
points (Theorem~\ref{thm:cub-transfer} in \S\ref{sec:rtc-transfer}); the avoidance of the line $\{Z_0=0\}$ is
Proposition~\ref{prop:rtc-avoid}.

\subsection{The packing}

The balls of capacities $\rho_i$ have total volume $\frac12\sum_i\rho_i^2$. The integrated profile is
$\frac12-\frac12(1-3h)^2$ (Lemma~\ref{lem:cub-volume}). We enlarge each $\rho_i$ by a slack $\delta$ to an auxiliary capacity
$a_i$ and choose $h$ so that the scalar budget after subtracting the enlarged ball capacities stays positive;
this budget is $\int_0^h\Pb$. As $\sum_i\rho_i^2\uparrow1$, the slack $\delta$, the defect $\frac13-h$ and the mean of $\Pb$
tend to zero.

\begin{lemma}[Auxiliary capacities]\label{lem:cub-auxiliary}
\lean{SymplecticPacking.Branch2.delta, SymplecticPacking.Branch2.aux, SymplecticPacking.Branch2.zeta,
SymplecticPacking.Branch2.height, SymplecticPacking.Branch2.tau, SymplecticPacking.Branch2.Pb,
SymplecticPacking.Branch2.Hyp.rho_lt_aux, SymplecticPacking.Branch2.Hyp.aux_lt_height,
SymplecticPacking.Branch2.Hyp.auxMax_lt_tau_mul, SymplecticPacking.Branch2.Hyp.tau_mul_lt_height,
SymplecticPacking.Branch2.Hyp.height_lt_third, SymplecticPacking.Branch2.Hyp.tau_pos, SymplecticPacking.Branch2.Hyp.tau_lt_one,
SymplecticPacking.Branch2.Hyp.sum_aux_sq_add_lt, SymplecticPacking.Branch2.Hyp.Pb_concave,
SymplecticPacking.Branch2.Hyp.Pb_continuous, SymplecticPacking.Branch2.Hyp.integral_Pb, SymplecticPacking.Branch2.Hyp.Pb_outer}
\uses{def:cub-model, lem:cub-volume}
Let $k\ge1$ and $0<\rho_1,\dots,\rho_k<\frac13$ with $\sum_i\rho_i^2<1$. Put $S_\rho=\sum_i\rho_i$, $D=1-\sum_i\rho_i^2$,
$\varpi=\frac13-\max_i\rho_i$,
\[
 \delta=\min\Big(\frac{\varpi}8,\ \frac{D}{8(S_\rho+k)},\ \frac18\Big),\qquad a_i=\rho_i+\delta,\qquad
 D_a=1-\sum_ia_i^2,
\]
\[
 \varsigma=\min\Big(\frac{\frac13-a_{\max}}2,\ \frac{D_a}{12}\Big),\qquad h=\frac13-\varsigma,\qquad
 \tau_\Delta=\frac{a_{\max}+h}{2h},
\]
with $a_{\max}=\max_ia_i$, and let $\Pb(p)=\Hb(p)-\sum_i(a_i-p)_+$ on $\Delta=[0,h]$. Then:
\begin{enumerate}[label=(\roman*)]
\item $\rho_i<a_i\le a_{\max}<\tau_\Delta h<h<\frac13$, and $0<\tau_\Delta<1$;
\item $\sum_ia_i^2+(1-3h)^2<1$;
\item $\Pb$ is continuous, concave and piecewise affine with breakpoints in $\{a_1,\dots,a_k\}$, and
\[
 \int_0^h\Pb(p)\,dp=\frac{1-\sum_ia_i^2-(1-3h)^2}{2}>0 ;
\]
\item $\Pb(p)=\Hb(p)\ge3(1-3h)>0$ for $p\in[\tau_\Delta h,h]$.
\end{enumerate}
If the $\rho_i$ are rational, so are $\delta$, $a_i$, $\varsigma$, $h$, $\tau_\Delta$ and the breakpoints and values of $\Pb$.
\end{lemma}

\begin{proof}
(ii) and positivity of $D_a$. Since $\delta\le\frac18$,
$\sum_ia_i^2-\sum_i\rho_i^2=2S_\rho\delta+k\delta^2\le\delta(2S_\rho+k)\le2\delta(S_\rho+k)\le D/4$, so $D_a\ge\frac34D>0$; also
$D_a\le1$. Then $(1-3h)^2=9\varsigma^2\le9D_a^2/144=D_a^2/16<D_a$, i.e.\ $\sum_ia_i^2+(1-3h)^2<1$.

(i) $\delta>0$ gives $a_i>\rho_i$, and $a_{\max}\le\max_i\rho_i+\varpi/8<\frac13$, so $\varsigma>0$ and $h<\frac13$. Further
$h\ge\frac13-\frac12(\frac13-a_{\max})=\frac12(\frac13+a_{\max})>a_{\max}$. Hence $\tau_\Delta h=\frac12(a_{\max}+h)$ lies strictly
between $a_{\max}$ and $h$, and $0<\tau_\Delta<1$.

(iii) $\Hb$ is affine and each $-(a_i-p)_+$ is concave and piecewise affine with the single breakpoint $a_i$. Since
$a_i<h$, $\int_0^h(a_i-p)_+\,dp=a_i^2/2$, and $\int_0^h\Hb=3h-\frac92h^2=\frac12\big(1-(1-3h)^2\big)$ by
Lemma~\ref{lem:cub-volume}. Positivity is (ii).

(iv) For $p\ge\tau_\Delta h>a_{\max}$ every $(a_i-p)_+$ vanishes, and $\Hb(p)=3(1-3p)\ge3(1-3h)$.

Rationality is clear from the formulas.
\end{proof}

With these data the packing in the model and the transfer give the four-dimensional embeddings,
Theorem~\ref{thm:cub-packing} in \S\ref{sec:rtc-packing}; the parameters of the equal-capacity branch are in
Remark~\ref{rem:rtc-equal}.
